\section{Restrições de Integridade}
\label{sec:integridade}

As regras de formato e de valor são garantidas pelo próprio banco, e não só
pela API. Assim, um dado inválido não entra nem por um script, nem por uma
consulta manual, nem por um bug.

\subsection{Chaves estrangeiras}

Toda chave estrangeira tem uma ação definida para quando a linha referenciada
é apagada:

\begin{itemize}
  \item \campo{ON DELETE CASCADE} nas tabelas associativas e nas que dependem
    de um dono: apagar uma playlist apaga os vínculos dela com músicas e
    membros. Na prática, isso quase nunca acontece, porque a aplicação usa
    remoção lógica.
  \item \campo{ON DELETE RESTRICT} em \campo{plan.type\_id}: um tipo de plano
    em uso não pode ser apagado.
\end{itemize}

\subsection{Restrições \texttt{CHECK}}

Cada \campo{CHECK} recusa a linha que viola a regra, com um erro que a API
traduz para a pessoa. A \cref{tab:checks} lista as principais.

\begin{table}[htbp]
  \centering
  \caption{Principais restrições \texttt{CHECK}}
  \label{tab:checks}
  \small
  \begin{tabularx}{\textwidth}{@{}L{0.9}L{1.1}@{}}
    \toprule
    \poliCabTabela{Restrição} & \poliCabTabela{Regra} \\
    \midrule
    \campo{rating BETWEEN 1 AND 5} & Nota da avaliação e média da música \\
    \campo{duration > 0} & Toda música tem duração positiva, em segundos \\
    \campo{price >= 0}, \campo{user\_amount > 0} & Preço e limite de pessoas do tipo de plano \\
    \campo{color \textasciitilde{} '\^{}\#[0-9A-Fa-f]\{6\}\$'} & Cores de gênero e de playlist no formato \campo{\#RRGGBB} \\
    \campo{email = lower(email)} & Emails de contas e de convites sempre em minúsculas \\
    \campo{role IN ('admin', 'assinante')} & Papéis possíveis de uma conta \\
    \campo{year >= 1900} & Ano do álbum; o limite ``até o ano atual'' fica no domínio, porque muda com o tempo \\
    \campo{track\_number > 0}, \campo{position >= 0} & Número da faixa no álbum e posição na playlist \\
    \campo{music\_id IS NULL OR event\_type IN (...)} & Só eventos de reprodução (\campo{music\_play}, \campo{music\_skip}, \campo{music\_listened}) apontam para uma música \\
    \campo{CHECK (singleton)} & \campo{system\_data} tem uma única linha \\
    \bottomrule
  \end{tabularx}
\end{table}

\FloatBarrier
\subsection{Unicidade}

Várias regras de unicidade valem só entre os registros ativos. Elas usam
\textbf{índices únicos parciais}, com \campo{WHERE deleted\_at IS NULL}, para
que um registro removido não bloqueie um novo:

\begin{itemize}
  \item nome de gênero único entre os ativos;
  \item nome de artista único entre os ativos, sem diferenciar maiúsculas de
    minúsculas (índice sobre \campo{lower(name)});
  \item uma avaliação ativa por pessoa e música.
\end{itemize}

Outras são \campo{UNIQUE} comuns: o email da conta (que continua reservado
mesmo depois da remoção), o token da sessão, o hash do token do convite, a
playlist de curtidas de cada conta e a denúncia de uma pessoa para um
comentário.

Duas unicidades de posição merecem destaque. O número da faixa é único dentro
do álbum (\campo{UNIQUE (album\_id, track\_number)}), e a posição é única
dentro da playlist. Na playlist, a constraint é \campo{DEFERRABLE INITIALLY
DEFERRED}: reordenar as músicas troca as posições num único \campo{UPDATE}, e
duas músicas passam um instante pela mesma posição. Com a verificação adiada,
o Postgres só confere a regra no \campo{COMMIT}.

\begin{listing}[htbp]
  \caption{Tabela de convites do plano familiar
    (\campo{V261105\_151000\_\_Adiciona\_convites\_do\_plano.sql})}
  \label{cod:convite}
  \begin{minted}{sql}
CREATE TABLE plan_invitation (
    id          integer      GENERATED BY DEFAULT AS IDENTITY,
    plan_id     integer      NOT NULL,
    email       varchar(254) NOT NULL,
    token_hash  varchar(64)  NOT NULL,
    expires_at  timestamptz  NOT NULL,
    accepted_at timestamptz,
    canceled_at timestamptz,
    created_at  timestamptz  NOT NULL DEFAULT now(),
    CONSTRAINT plan_invitation_pkey PRIMARY KEY (id),
    CONSTRAINT plan_invitation_token_hash_key UNIQUE (token_hash),
    CONSTRAINT plan_invitation_email_lower_check CHECK (email = lower(email)),
    CONSTRAINT plan_invitation_plan_id_fkey FOREIGN KEY (plan_id)
        REFERENCES plan (id) ON DELETE CASCADE
);
  \end{minted}
\end{listing}

O \cref{cod:convite} reúne os elementos desta seção: chave primária gerada,
chave estrangeira com ação definida, \campo{UNIQUE} e \campo{CHECK}, todos com
nome no padrão do projeto. O estado do convite não é uma coluna: ele está
pendente enquanto \campo{accepted\_at} e \campo{canceled\_at} são nulos e
\campo{expires\_at} está no futuro.

\FloatBarrier
\subsection{Regras que ficam no domínio}

Algumas regras dependem de contar linhas ou do momento atual, e um
\campo{CHECK} não as expressa. Elas são verificadas nos casos de uso da API:

\begin{itemize}
  \item todo plano tem ao menos uma pessoa, e o Família não passa de quatro,
    contando os convites pendentes;
  \item toda música e todo álbum têm ao menos um artista;
  \item toda playlist tem ao menos um membro;
  \item o ano do álbum não passa do ano atual;
  \item uma música só conta como ouvida com mais de 30 segundos de reprodução
    (\campo{LISTENED\_THRESHOLD\_SECONDS}).
\end{itemize}
